% Nota's LaTeX render, exercised. Press ⇧⌘P to cycle editor → split →
% reading, or ⌥⇧⌘P to jump straight to reading. This whole line is a comment
% and renders nothing; an escaped \% is a percent sign.
\documentclass[11pt]{article}
\usepackage{amsmath}
\usepackage{graphicx}
\newcommand{\nota}{\textbf{Nota}}

\title{Reading \LaTeX{} in Nota}
\author{The Nota preview}
\date{September 2026}

\begin{document}
\maketitle

\begin{abstract}
Nota reads an article-level subset of \LaTeX{} and lays it out to be read ---
not to be typeset. The preamble above renders nothing at all; only the title,
the author and the date were taken out of it.
\end{abstract}

\section{Text}
\label{sec:text}

Paragraphs are parted by a blank line, and a source line break is just a space,
so this sentence
runs on.

You can set text \textbf{bold}, \textit{italic}, \emph{emphasised},
\texttt{monospaced}, \underline{underlined} and \textsc{small caps}, and nest
them: \textbf{\textit{both at once}}.

Punctuation is TeX's: `single quotes', ``double quotes'', an en--dash, an
em---dash, a non~breaking space, and an ellipsis\ldots{} A hard break comes
from a double backslash,\\ like that.

A footnote\footnote{Which shows as a bracketed note in the secondary colour,
right where the mark stood.} sits in the running text. So do citations
\cite{knuth1984} and cross-references to Section \ref{sec:maths} and to
Equation \eqref{eq:gauss}. Links are clickable: \url{https://example.com} and
\href{https://example.com}{a named link}.

\subsection{Lists}

\begin{itemize}
  \item Bullets, from \texttt{itemize}.
  \item They nest:
  \begin{enumerate}
    \item Numbered, from \texttt{enumerate}.
    \item And the numbers count.
  \end{enumerate}
  \item Back out at the first level.
\end{itemize}

\begin{description}
  \item[Block sync] Clicking a block here puts the caret on its source.
  \item[Both ways] Moving the caret over there lights the block here.
\end{description}

\subsubsection{Quotations}

\begin{quote}
A quoted paragraph takes a leading hairline bar, the same one a Markdown block
quote gets.
\end{quote}

\section{Maths}
\label{sec:maths}

Inline maths sits in the sentence: the roots are
$x = \frac{-b \pm \sqrt{b^2 - 4ac}}{2a}$, Euler's identity is
$e^{i\pi} + 1 = 0$, and \(a^2 + b^2 = c^2\) is Pythagoras.

An unnumbered display formula stands on its own:
\[
  \lim_{x \to 0} \frac{\sin x}{x} = 1
\]

A numbered one carries its number at the trailing edge:

\begin{equation}
\label{eq:gauss}
\int_{-\infty}^{\infty} e^{-x^2}\,dx = \sqrt{\pi}
\end{equation}

An \texttt{align} splits on the double backslash, drops the alignment
ampersands, and numbers every row in its own right:

\begin{align}
(a + b)^2 &= a^2 + 2ab + b^2 \\
(a - b)^2 &= a^2 - 2ab + b^2
\end{align}

A starred environment is unnumbered:

\begin{gather*}
\sum_{i=1}^{n} i = \frac{n(n+1)}{2}
\end{gather*}

\section{Blocks}

\subsection{Code}

Verbatim is untouched --- the dollars and backslashes in here are characters,
not delimiters:

\begin{verbatim}
for i in 1..$n; do
    echo "$i \frac{1}{2}"
done
\end{verbatim}

\subsection{Tables}

\begin{table}
\begin{tabular}{lrl}
\hline
Environment & Rows & Numbered \\
\hline
equation & 1 & yes \\
align & many & every row \\
gather* & many & no \\
\hline
\end{tabular}
\caption{What the maths environments do.}
\end{table}

\subsection{Pictures}

An \verb-\includegraphics- whose path resolves beside this file draws the
picture; one that resolves to nothing shows the path instead, as this one does:

\begin{figure}
\includegraphics[width=0.6\textwidth]{diagram}
\caption{A picture that is not there.}
\end{figure}

\section{What the subset has never heard of}

A command it does not know renders as its own raw source, inline in the
code-span style: \marginpar{like this one}. So does a use of the
\verb-\newcommand- from the preamble: \nota.

An environment it does not know renders whole, as a code block:

\begin{theorem}
Every unknown environment comes back as itself. Never a blank, never a broken
page --- the same rule the maths engine keeps.
\end{theorem}

\begin{center}
That is the whole of it.
\end{center}

\end{document}
